\section{Categories}

\begin{definition}
    A category $\mathcal{C}$ consists of:
    \begin{itemize}
        \item A class of objects $\operatorname{ob}(\mathcal{C})$.
        \item A set of morphisms $\operatorname{Hom}_{\mathcal{C}}(A, B)$ for each pair of objects $A, B \in \mathcal{C}$.
        \item A composition map $\operatorname{Hom}_{\mathcal{C}}(B, C) \times \operatorname{Hom}_{\mathcal{C}}(A, B) \to \operatorname{Hom}_{\mathcal{C}}(A, C)$, $(\phi, \psi) \mapsto \phi \circ \psi$ for each triple of objects $A, B, C \in \mathcal{C}$.
    \end{itemize}
\end{definition}

It is assumed that objects and morphisms satisfy the following conditions:
\begin{itemize}
    \item If $(A,B)\neq (C,D)$, then
          \[
              \operatorname{Hom}(A,B)\cap \operatorname{Hom}(C,D)=\varnothing.
          \]

    \item If
          \[
              f\in \operatorname{Hom}(A,B), \qquad
              g\in \operatorname{Hom}(B,C), \qquad
              h\in \operatorname{Hom}(C,D),
          \]
          then
          \[
              (h\circ g)\circ f = h\circ (g\circ f).
          \]

    \item (Unit) For every object $A$ we have an element
          \[
              1_A \in \operatorname{Hom}(A,A)
          \]
          such that
          \[
              f\circ 1_A = f \qquad \text{for any } f\in \operatorname{Hom}(A,B),
          \]
          and
          \[
              1_A\circ g = g \qquad \text{for any } g\in \operatorname{Hom}(B,A).
          \]
\end{itemize}

\begin{example}
    The following are examples of categories.
    \begin{enumerate}
        \item The category of sets: the objects are sets, and the morphisms are maps. It is denoted by $\mathbf{Set}$.

        \item The category of groups: the objects are groups, and the morphisms are group homomorphisms. It is denoted by $\mathbf{Grp}$.

        \item The category of abelian groups: the objects are abelian groups, and the morphisms are group homomorphisms. It is denoted by $\mathbf{Ab}$.

        \item The category of $R$-modules: the objects are $R$-modules, and the morphisms are $R$-module homomorphisms. It is denoted by $R\text{-}\mathbf{Mod}$(left $R$-modules, for right $R$-modules, we use $\mathbf{Mod}-R$).

        \item The category of finitely generated $R$-modules: the objects are finitely generated $R$-modules, and the morphisms are $R$-module homomorphisms. It is denoted by $R\text{-}\mathbf{mod}$.

        \item The category of topological spaces: the objects are topological spaces, and the morphisms are continuous maps. It is denoted by $\mathbf{Top}$.
    \end{enumerate}
\end{example}

\begin{definition}[Subcategory]
    A category $\mathcal{D}$ is called a \emph{subcategory} of a category $\mathcal{C}$ if
    \begin{enumerate}
        \item the objects of $\mathcal{D}$ form a subclass of the objects of $\mathcal{C}$;
        \item for any objects $A,B$ of $\mathcal{D}$, we have
              \begin{itemize}
                  \item $\operatorname{Hom}_{\mathcal{D}}(A,B)\subseteq \operatorname{Hom}_{\mathcal{C}}(A,B)$;
                  \item the identity morphisms and the composition of morphisms in $\mathcal{D}$ are the same as those in $\mathcal{C}$.
              \end{itemize}
    \end{enumerate}

    The subcategory $\mathcal{D}$ of $\mathcal{C}$ is called \emph{full} if
    \[
        \operatorname{Hom}_{\mathcal{D}}(A,B)=\operatorname{Hom}_{\mathcal{C}}(A,B)
    \]
    for any objects $A,B$ of $\mathcal{D}$.
\end{definition}

\begin{example}
    The category $\mathbf{Ab}$ is a full subcategory of $\mathbf{Grp}$.

    The category $R\text{-}\mathbf{mod}$ is a full subcategory of $R\text{-}\mathbf{Mod}$.

    The category $\mathbf{Grp}$ is not a full subcategory of $\mathbf{Set}$.
\end{example}

\begin{definition}
    Let $\mathcal{C}$ and $\mathcal{D}$ be categories. A \emph{covariant functor}
    \[
        F:\mathcal{C}\to\mathcal{D}
    \]
    consists of the following data:
    \begin{itemize}
        \item for every object $A$ of $\mathcal{C}$, an object $F(A)$ of $\mathcal{D}$;
        \item for each pair of objects $A,B$ of $\mathcal{C}$, a map
              \[
                  F:\operatorname{Hom}_{\mathcal{C}}(A,B)\to \operatorname{Hom}_{\mathcal{D}}(F(A),F(B)),
                  \qquad
                  f\mapsto F(f).
              \]
    \end{itemize}
    These data are required to satisfy the following conditions:
    \begin{enumerate}
        \item If $g\circ f$ is defined in $\mathcal{C}$, then
              \[
                  F(g)\circ F(f)=F(g\circ f).
              \]
        \item For every object $A$ of $\mathcal{C}$,
              \[
                  F(1_A)=1_{F(A)}.
              \]
    \end{enumerate}
\end{definition}

\begin{proposition}
    Functors preserve diagram:
    \begin{center}
        \begin{tikzcd}
            A \arrow[r, "f"] \arrow[rd, "g"'] & B \arrow[d, "h"] & {} \arrow[r, "F"] & {} & F(A) \arrow[r, "F(f)"] \arrow[rd, "F(g)"'] & F(B) \arrow[d, "F(h)"] \\
            & C                &                   &    &                                           & F(AC
        \end{tikzcd}
    \end{center}
\end{proposition}

Let $M$ and $N$ be $R$-modules. Then
\[
    \operatorname{Hom}_R(M,N)
    =
    \{\, f:M\to N \mid f \text{ is an } R\text{-module homomorphism} \,\}.
\]
This is an abelian group.

Moreover,
\[
    \operatorname{End}_R(M)=\operatorname{Hom}_R(M,M),
\]
and $\operatorname{End}_R(M)$ is a ring.

\begin{example}
    Let $M$ be an $R$-module. Then
    \[
        \operatorname{Hom}_R(M,-): R\text{-}\mathrm{Mod} \to \mathrm{Ab}
    \]
    is a functor.

    More precisely, it acts as follows:

    \[
        N \longmapsto \operatorname{Hom}_R(M,N),
        \qquad
        f:N \to N' \longmapsto \operatorname{Hom}_R(M,f).
    \]

    Given an $R$-module homomorphism
    \[
        f:N \to N',
    \]
    the induced map
    \[
        \operatorname{Hom}_R(M,f): \operatorname{Hom}_R(M,N) \to \operatorname{Hom}_R(M,N')
    \]
    is defined by
    \[
        \operatorname{Hom}_R(M,f)(h)=f\circ h
        \qquad
        (h\in \operatorname{Hom}_R(M,N)).
    \]

    This corresponds to the commutative diagram
    \[
        \begin{tikzcd}
            & M \arrow[dl, "h"'] \arrow[dr, "f\circ h"] & \\
            N \arrow[rr, "f"'] & & N'
        \end{tikzcd}
    \]

    Hence
    \[
        \operatorname{Hom}_R(M,-)
    \]
    is a covariant functor.
\end{example}

\begin{definition}
    Let $\mathcal C$ be a category. The \emph{opposite category} (or \emph{dual category}) $\mathcal C^{\mathrm{op}}$ is defined as follows:
    \begin{itemize}
        \item $\operatorname{Ob}(\mathcal C^{\mathrm{op}})=\operatorname{Ob}(\mathcal C)$;
        \item for any objects $A,B$,
              \[
                  \operatorname{Hom}_{\mathcal C^{\mathrm{op}}}(A,B)
                  =
                  \operatorname{Hom}_{\mathcal C}(B,A).
              \]
    \end{itemize}
    The composition in $\mathcal C^{\mathrm{op}}$ is defined by reversing the order of composition in $\mathcal C$: if
    \[
        f \in \operatorname{Hom}_{\mathcal C^{\mathrm{op}}}(A,B),
        \qquad
        g \in \operatorname{Hom}_{\mathcal C^{\mathrm{op}}}(B,C),
    \]
    then
    \[
        g\circ f \text{ in } \mathcal C^{\mathrm{op}}
    \]
    is defined to be
    \[
        f\circ g \text{ in } \mathcal C.
    \]
    We call $\mathcal C^{\mathrm{op}}$ the dual category of $\mathcal C$.
\end{definition}

This is illustrated by the following diagrams:
\[
    \begin{tikzcd}[column sep=large,row sep=large]
        A \arrow[r,"f"] \arrow[dr,"g\circ f"']& B \arrow[d,"g"] \\
        & C
    \end{tikzcd}
    \qquad\text{in } \mathcal C^{\mathrm{op}},
\]
while the corresponding picture in $\mathcal C$ is
\[
    \begin{tikzcd}[column sep=large,row sep=large]
        A & B \arrow[l,"f"'] \\
        & C \arrow[ul,"f\circ g"] \arrow[u,"g"']
    \end{tikzcd}
    \qquad\text{in } \mathcal C.
\]

\begin{definition}
    A \emph{contravariant functor} from a category $\mathcal C$ to a category $\mathcal D$ is a covariant functor
    \[
        \mathcal C^{\mathrm{op}} \to \mathcal D.
    \]
\end{definition}

\begin{example}
    Let $M$ be an $R$-module. Then
    \[
        \operatorname{Hom}_R(-,M): R\text{-}\mathrm{Mod} \to \mathrm{Ab}
    \]
    is a contravariant functor.
\end{example}

More precisely, it is defined on objects by
\[
    N \longmapsto \operatorname{Hom}_R(N,M),
\]
and on morphisms by sending an $R$-module homomorphism
\[
    f:N \to N'
\]
to the induced group homomorphism
\[
    \operatorname{Hom}_R(f,M): \operatorname{Hom}_R(N',M) \to \operatorname{Hom}_R(N,M),
\]
given by
\[
    \operatorname{Hom}_R(f,M)(h)=h\circ f
    \qquad
    (h\in \operatorname{Hom}_R(N',M)).
\]

This corresponds to the commutative diagram
\[
    \begin{tikzcd}
        N \arrow[d,"f"'] \arrow[dr,"h\circ f"] & \\
        N' \arrow[r,"h"'] & M
    \end{tikzcd}
\]

